\section{Data for Diagnostic Module} \label{sec2.12}

\textbf{General remarks} 

\medskip 

The data in this block are used to define a line-of-sight (LOS) across the
computational domain, along which line integrals $I$
\begin{equation}
  \label{eqLOS} I = \int_{LOS} g(l) dl = \int_{P_1}^{P_2} g(l) dl
\end{equation}
are evaluated (by launching a ``virtual" test particle along the LOS and
scoring the response function $g(x)$ along its track, using the EIRENE particle
tracing routines.
%

This is done in subroutine LININT, which is called from subroutine DIAGNO at
the end (post processing phase) of an EIRENE run.
The spatial dependence of the function $g(l)$ is defined as function of one or
more of the estimated volume averaged tallies (e.g., atom density), and/or
input tallies (e.g.\ plasma temperatures) and/or a small set of preprogrammed
additional tallies, such as volumetric line emissivities for various H and He
lines (e.g.\ Ba\_alpha.f for the Balmer alpha emission rate per volume, etc.).
At present there are two (version 2004 and older) or three preprogrammed
functions $g(l)$, and the option to call a user supplied integrand.

Firstly there are a number of Lyman- and Balmer series volume source rates
(emissivity), see Subr.\ SIGLINE.
Each hydrogenic line emissivity is available for coupling to up to 6 states
(``donor states") and linear in the corresponding density (output tally),
namely to the (ground) states:
$H, H_2, H_2^+, H^+, H^-$ and $H_3^+$	(and their isotopes,
isotopomeres\dots).

Secondly the neutral atom charge-exchange source rate can be integrated along a
line of sight for a given energy of the impacting plasma ion (Subr.\ SIGCX).
This routine includes re-absorption along the line of sight.
A LOS spectrum (also: ``side-on spectrum''), energy resolved with up to NCHEN
energies, may be obtained by this procedure.

Thirdly (version 2005 and younger) the side-on radiances of selected lines
(photon test particle species), including reabsorption, reflection, line
broadening, etc., can be obtained (Subr.\ SIGRAD).
A number of different emission line shape profiles is available for these.

For testing and \ordinalnum{3} party defined further integrands $g(x)$ along
lines of side there is a user defined option (Subr.\ SIGTST).

In 2018 a more general option for definition of emission lines (in NCHTAL = 2
option) was introduced. It allows the definition of emission lines via the 
Eirene input file. The population coefficients, QSS ratios, etc. needed for
setting up volumetric line emissivity profiles as further additional tally ADDV
need to be specified in input block 4. The ADDV tallies may then be used 
for line of sight integration, along the chords defined further below.

Distinct from the other addv tallies from block 10a these
further addv tallies (beyond NADVI) are currently apparently neither
scaled, averaged, integrated, nor do they have text (units, species) assigned.

One of the keywords ``DEFINE\_LINES'' or ``DEFAULT\_LINES'' needs to specified 
in the input to determine if the originally implemented Lyman- and 
Balmer series volume source rates should be used or line definitions are 
given. Please be aware: there is no mixing of default and 
selfdefined lines. If at least one selfdefined line is to be used all 
emission lines need to be defined in the input.



\medskip

\textbf{The Input Block} 

\begin{lstlisting}
*** 12. Data for Diagnostic Module
      KEYWORD
      IF (KEYWORD == ``DEFINE_LINES'') THEN
        NUM_LINES, MOD_ADDV
        DO ILINE=1, NUM_LINES
**
          LINE_NAME
          NUM_COMPO
          EINSTEIN TRANS_EN
          DO JCOMP=1,NUM_COMPO
            COMPO_NAME
            NUM_CONTRIB, IRC
            DO KCONTR = 1, NUM_CONTRIB
              ISP, ITP, IRATIO
              IF (IRATIO > 0) THEN
                IRC_RAT(1)
                IF (IRATIO == 2) THEN
                  IRC_RAT(2), (ISP_RAT(I), ITP_RAT(I) I=1, 2)
                ENDIF
              ENDIF
            ENDDO
          ENDDO
        ENDDO
      ENDIF 
      NCHORI  NCHENI
      IF (NCHORI.GT.0) THEN
        DO 121 ICHORI=1,NCHORI
          TXTSIG
          NSPTAL NSPSCL NSPNEW NSPCHR
          NSPSTR NSPSPZ NSPINI NSPEND NSPBLC NSPADD
          EMIN1  EMAX1 ESHIFT
          IPIVOT  XPIVOT  YPIVOT  ZPIVOT
          ICHORD  XCHORD  YCHORD  ZCHORD
  121   CONTINUE
        PLCHOR PLSPEC PRSPEC PLARGL PRARGL
      ENDIF
\end{lstlisting} 

\medskip 

\textbf{Meaning of the Input Variables for Diagnostic Module}

\medskip 

\begin{description}
  \item[KEYWORD] Keyword indicating if the originally implemented Lyman- and 
        Balmer series volume source rates should be used or 
        line definitions are provide. 
        \begin{description}
          \item[$=$ DEFAULT\_LINES] the implemented Lyman- and 
               Balmer series volume source rates are used
          \item[$=$ DEFINE\_LINES] definition of emission lines is expected
        \end{description}
  \item[NUM\_LINES] Total number of different emission lines
  \item[MOD\_ADDV] flag for memory management of ADDV tallies
        \begin{description}
          \item[$= 0$] memory saving mode
               Only the components of the latest calculated emission line
               are kept in storage. A new calculated emission line 
               overwrites the previously calculated data. If a previously
               calculated emission line is used again for a chord lateron
               it needs to be recalculated.
          \item[$= 1$] All components of all emission lines are stored
               and kept. This option is much more storage consuming but
               also faster than MOD\_ADDV=0.
        \end{description}
   \item[LINE\_NAME] name of the emission line, e.g. BA\_ALPHA.

         LINE\_NAME can be used to identify the particular emission
         profile in the specification of a line of sight.
   \item[NUM\_COMPO] Number of components of the line
   \item[EINSTEIN] radiative transition rate (1/sec)
   \item[TRANS\_EN] transition energy, identifyer for line
   \item[COMPO\_NAME] name of the component, 

         used for identification of the component in the emissivity output
   \item[NUM\_CONTRIB] Number of contributions to the component of the line
   \item[IRC] Number of a reaction defined in input block 4 for a
         population coefficient such as H.12 2.2.5a
   \item[ISP] species index of density which is used for building the 
         emission line
   \item[ITP] particle type
        \begin{description}
          \item[$= 0$] photon density PDENPH is used
          \item[$= 1$] atom density PDENA is used
          \item[$= 2$] molecule density PDENM is used
          \item[$= 3$] ion density PDENI is used
          \item[$= 4$] bulk ion density DIIN is used
          \item[$= 5$] electron density DEIN is used
        \end{description}
   \item[IRATIO] number of density ratios to be used for this contribution,
         maximal 2 ratios are allowed.

         The use of a ratio is required if the true parent density is not
         available (or in QSS mode). Then the first ratio converts from 
         density(1) the to ``true'' density for this component.

         density(1) is taken as "intermediate" parent density. Fetch 
         reduced population coefficent and density ratio  
         ratio1="density"/"density(1)" will be applied,
         to turn density(1) into "density"

         e.g. density    = H2+

         \hspace{0.6 true cm} density(1) = H2

         \hspace{0.6 true cm} ratio1     = [H2+]/[H2]

         This works when the second species involved in loss and gain rates
         for species H2+ from H2 is the same
   \item[IRC\_RAT] Number of a reaction defined in input block 4 for providing
         a density ratio
   \item[ISP\_RAT] species number of the density to be used with ratio I, I=1,2
   \item[ITP\_RAT] type of particle density to be used with ratio I, I=1,2
   \item[NCHORI] Total number of different line of sights
   \item[NCHENI] $\mid$NCHENI$\mid$ is the total number of energies, at which the
        spectrum is evaluated (irrelevant e.g.\ for total signals, such as Lyman and
        Balmer emissivity without line shape resolution).

        The energy grid is equidistant on a linear scale, if NCHENI $\ge$ 0, and
        equidistant on a logarithmic scale otherwise.
        (Same as for NSPS parameter in input block 10F for spectrally resolved default
        tallies.)
  \item[TXTSIG] Text in printout at the beginning of the data from this line of
        sight integral.
  \item[NSPTAL] Flag for choice of preprogrammed function which is ``line
        integrated".
        \begin{description}
          \item[$= 1$] charge-exchange source rate (SIGCX)
          \item[$= 2$] Hydrogen line
                emission source rate (SIGLINE)
          \item[$= 3$] spectral radiance of photonic lines
                (SIGRAD) (new in versions 2005 and younger)
          \item[$= 10$] user supplied
                integrand (SIGUSR) 

                (this was option NSPTAL =3, in versions 2004 and older)
          \item[$= 11$] integrand is line of sight (SIGLOS) 
          \item[$= 12$] back ground (input) tallies
                user supplied integrand (SIGPLA) 
          \item[$= 13$] volume averaged (output) tallies
                user supplied integrand (SIGEIR) 
        \end{description}
  \item[NSPSCL] Flag for choice of linear or logarithmic axes in plots of spectra
        (vs.\ energy or wavelength, resp.) and of source term distribution along 
        line of sight.
        \begin{description}
          \item[$= 0$] both axes linear
          \item[$= 1$] x axis linear, y axis logarithmic
          \item[$= 2$] x axis logarithmic, y axis linear
          \item[$= 3$] both axes
                logarithmic
        \end{description}
  \item[NSPNEW] Flag for the choice whether a
        spectrum is plotted on the same picture as the previous one (NSPNEW = 0) or
        onto a new graph (otherwise).
  \item[NSPCHR] (Proprietary option, not for \ordinalnum{3} party application) If
        NSPCHR.gt.0, then automatically all grid cells of the computational volume
        along the chord are identified and marked in pre-processing.
        During the particle tracing phase then, energy resolved spectra (input block
        10f) are computed (scored) in all these cells directly from Monte Carlo
        histories, by automatically augmenting input block 10f (additional energy
        resolved tallies: ''spectra") correspondingly.
        These spectra are line-of-sight spectra in the direction of the chord specified
        here (direction SPCVX, SPCVY, SPCVZ in augmented input block 10f is taken to be
        a unit vector along this present line of sight).
  \item[NSPSTR] Index for stratum, which is to be used for line integration
        (NSPSTR = 0: sum over strata).
  \item[NSPSPZ] Index for species which is to be used for line integration

        In case NSPTAL=1:

        Index for atomic species IATM, with IATM $\leq$ NATM, for which charge-exchange
        spectrum is to be computed (NSPSPZ = 0: sum over atom species index)

        In case NSPTAL=2:

        (NSPSTR to be written, current default, hard wired: hydrogenic atoms,
        molecules, test ions, background ions).

        Number of contribution to line intensity, as programmed in Subr.\
        Ba\textsubscript{Alpha}, Ba\textsubscript{Gamma}, Ly\textsubscript{Beta}, etc.
        (Currently up to 6 contributions for each H atom spectral line, and the total)
        \begin{description}
          \item[$= 1$] coupling to ground state atom: $H(1s)$
          \item[$= 2$] coupling to continuum, bulk ion: $H^+$
          \item[$= 3$] coupling to
                ground state molecule: $H_2$
          \item[$= 4$] coupling to ground state molecular
                ion : $H_2^+$
          \item[$= 5$] coupling to negative ion: $H^-$ (new in versions
                1996 and younger)
          \item[$= 6$] coupling to: $H_3^+$  (new in versions 2012 and
                younger)
          \item[$= 0$] total, sum over all contributions to a particular line
                (was =~6, in versions 2011 and older)
        \end{description}

        In case NSPTAL=3:

        Index for photon species (line), i.e.\ for IPHOT, with IPHOT $\leq$ NPHOT, for
        which side-on spectrum is to be computed (NSPSPZ = 0: sum over photon species
        index, not ready)
  \item[NSPINI, NSPEND] only for NSPTAL=1:

        Multipliers for the maximum ion temperature Ti$_{max}$ found along line of
        sight, for temperature fitting.
        The CX ion temperature	is fitted from the CX line of side spectrum in the
        interval [NSPINI $\times$ Ti$_{max}$, NSPEND $\times$ Ti$_{max}$ ]

  \item[NSPBLC] Standard mesh block number of \ordinalnum{2} point on line of
        sight.
  \item[NSPADD] Additional cell number of \ordinalnum{2} point on line of sight.

        If NSPADD = 0, then this \ordinalnum{2} point must lie in standard mesh block
        NSPBLC.

        If NSPADD $\ne$ 0, then the block number NSPBLC must be NSPBLC = NBMLT+1, i.e.\
        the second point on the line of sight is in the ``additional cell region".
  \item[EMIN1, EMAX1] ~ 

        for NSPTAL=1,3,10:

        minimum and maximum energy for spectral resolution, respectively.

        for NSPTAL=2:

        \textcolor[rgb]{1.00,0.00,0.00}{(DR, Oct.\ 2016 currently not
          available, use old option, see below)}

        Parameter to identify the particular hydrogen line.
        Emin1 and Emax1 are interpreted as integer principal quantum numbers $n,m$ of
        the lower and upper state for the transition.

        EMIN1 = 1, EMAX1 = 3:  Lyman-beta line 

        EMIN1 = 2, EMAX1 = 6:  Balmer-delta line 

        EMIN1 = 2, EMAX1 = 5:  Balmer-gamma line 

        EMIN1 = 2, EMAX1 = 4:  Balmer-beta line 

        EMIN1 = 2, EMAX1 = 3:  Balmer-alpha line 

        Old input version (still maintained for backward compatibility of input files:
        EMIN1 is an energy parameter to identify the particular hydrogen line and EMAX1
        is not used.
        EMIN1 is given in \si{\electronvolt}, by $Ry \times \left(1/n^2-1/m^2\right)$,
        with $Ry=\num{13.6}$ (\si{\electronvolt}).

        EMIN1 = 12.089: Lyman-beta line 

        EMIN1 = 3.0222: Balmer-delta line 

        EMIN1 = 2.8560: Balmer-gamma line 

        EMIN1 = 2.5500: Balmer-beta line 

        EMIN1 = 1.8889: Balmer-alpha line 

        Other side on line emissivities can be obtained using the user supplied line of
        side integral SIGUSR (option NSPTAL=10) and analogy to the preprogrammed
        options, as well as the internal EIRENE hydrogen atom collisional radiative
        routine \emph{H-COLRAD.F} to obtain the required reduced population
        coefficients for $H^*(n)$.
        For the preprogrammed options these latter coefficients for $n=2,3,4,5,6$ are
        stored in AMJUEL, section H12, see this web page under: EIRENE AMS data files.

  \item[ESHIFT] (for NSPTAL=1, 3, 10 options only) 

        energy shift for spectral resolution in printout, and plot
\end{description}

The line integration is carried out along a line defined by two points, a pivot
point $P_{pivot}$ and a second point $P_{scnd}$.
The second point must lie inside mesh block no.
NSPBLC.
Starting from this ``second point", and moving in the direction towards the
``first point"	    $P_{pivot}$, the first intersection $P_1$ with a non
transparent additional or non default standard surface is computed.
This is the starting point for line of sight integration.
Then the line integration is carried out starting a ``virtual particle" from
this point $P_1$, into the computational area, until the next intersection
$P_2$ with any non transparent additional or non default standard surface is
found.

\medskip

Note 1:

When specifying the coordinates of the points $P_{pivot}$ and $P_{scnd}$ to
select a particular line of sight, the same rules must be followed as also for
specifying birth point locations of ``real" particles (input block 7).
These starting points are assumed to be given \textbf{inside a cell, not 
exactly on a cell boundary}.
For example coordinates z0=0.0 must be avoided if a z-grid surface coincides
with this position.
This latter is often the case e.g.\ already when default options for toroidal
curvature (NLTRA option in block 2c) are used, even in otherwise only 1D
(radial, x) or only 2D (radial and poloidal, x,y) resolved cases.
The line of side points then must be specified at least with a small
incremental distance to such coordinate surfaces, e.g.\ z0=1.0E-5 in the
example.

Note 2:

The line integration (virtual particle flight) stops at the first
non-transparent surface encountered along the line of sight.
In cases when only the boundary plasma domain is simulated, the interface to
the core plasma region is often such a ``non-transparent surface" (mostly set
to be purely absorbing).
To allow line of sight \textbf{integration across this central core region}
e.g.\ to capture both  top and bottom machine contributions along a line of
sight, the following work-around can be used:

a) set the core boundary transparent, and switch (ILSWCH option block 3) into
an unused additional cell,
e.g.\ additional cell no.\ NC.

b) define a fixed ``virtual core plasma" in this additional cell no.\ NC, e.g.\
by the options in input block 8,
such that during ``real" particle simulations the test particles get either
absorbed in the core of reflected (by scattering)
back into the boundary plasma domain.

For example in EMC3-EIRENE runs high plasma density and temperature are used in
virtual inner core cells, which hence to not contribute significantly and
artificially to line of sight integrals of atomic or molecular emission lines,
which cross that core region.

\medskip

The next card specifies the first (``pivot") point.

\begin{description}
  \item[IPIVOT] (only needed for NLTRA option, ``toroidal approximation",
        sub-block 2c)
        \begin{description}
          \item[1 $\leq$ IPIVOT $\leq$ NTTRA-1]
                (currently no available, error exit) 

                number of local toroidal co-ordinate system (NTTRA: see sub-block 2c), in which
                this pivot point is specified.
                The pivot point is then given in Cartesian coordinates is this local system
          \item[IPIVOT = 0] ~ 

                pivot point is given in global cylindrical coordinates, 

                XPIVOT,YPIVOT,ZPIVOT = $r,z,\phi$, with $\phi$ in degrees.
                The corresponding toroidal block numbers ITTRA are found automatically.
        \end{description}
  \item[XPIVOT] \ordinalnum{1} co-ordinate of pivot point for line of sight,
        e.g.\ $x$
  \item[YPIVOT] \ordinalnum{2} co-ordinate of pivot point for line of
        sight, e.g.\ $y$
  \item[ZPIVOT] \ordinalnum{3} co-ordinate of pivot point for
        line of sight, e.g.\ $z$
\end{description}
The next card specifies the second
(``inside") point.

The second point on the line of sight must either be inside the first (radial
or x-) mesh, or the additional cell number of the mesh cell, to which this
point belongs, must be specified by the NSPBLC and NSPADD flags.

(For otherwise the initial point for the line integration cannot be found
automatically.)

\begin{description}
  \item[ICHORD] (only needed for NLTRA option, sub-block 2c) 

        Meaning analogous to that of first point flag IPIVOT:

  \item[XCHORD] \ordinalnum{1} co-ordinate of second point for line of sight
  \item[YCHORD] \ordinalnum{2} co-ordinate of second point for line of sight
  \item[ZCHORD] \ordinalnum{3} co-ordinate of second point for line of sight
  \item[PLCHOR] the lines of sight are plotted into (2D or 3D) geometry plots.
        This, however, is automatically turned off if other plots are done between
        geometry plots (initialization phase) and line-of-sight integration (post
        processing phase).
  \item[PLSPEC] the energy or wavelength resolved spectra, integrated along the
        lines of sight, are plotted (irrelevant in case of NSPTAL = 2).
  \item[PRSPEC] the energy or wavelength resolved spectra, integrated along the
        lines of sight, are printed (from subr.\ OUTSIG.f) (irrelevant in case of
        NSPTAL = 2).
  \item[PLARGL] the line of sight spectra, with their contributions spatially
        resolved along the lines of sight, are plotted.
  \item[PRARGL] the line of sight spectra, with their contributions spatially
        resolved along the lines of sight, are printed.
        (directly from subr.\ LININT.f)
\end{description}

\subsection{Line of sight: charge-exchange spectrum}
\label{sec2.12.1}

\subsection{Line of sight: line emissivity}\label{sec2.12.2}

\subsection{Line of sight: line shape}\label{sec2.12.3}

\subsection{Line of sight: user defined integral}\label{sec2.12.4}
